%% =====================================================================
%%  T-10-01  Value as a Function of Availability
%%  -------------------------------------------------------------------
%%  One idea per file. Core is mandatory. Slots in canonical order.
%%  Max 700 words. No layout commands. No filler. New source material
%%  for this entry goes in sources/<BOOK>/T-10-01.tex -- never by
%%  rewriting this file (docs/RULES.md R0.1).
%% =====================================================================
%% @kind: topic
%% @id: T-10-01
%% @title: Value as a Function of Availability
%% @subtitle: The rarer a thing appears, the more it is worth
%% @chapter: c10
%% @order: 10
%% @status: stable
%% @sources: B4
%% @updated: 2026-09-19

\topic{T-10-01}{Value as a Function of Availability}{The rarer a thing appears, the more it is worth}

\begin{core}
Opportunities look more valuable when they look less available. The judgement is about scarcity rather than about the thing, so scarcity can be produced and the value follows it without the object changing at all.
\end{core}
\begin{mechanism}
Scarcity carries information. A thing that is hard to obtain has usually survived some test --- of demand, of difficulty, of selection --- and treating rarity as a proxy for quality is normally a reasonable shortcut. It is also a shortcut about loss rather than about gain: as availability falls, what is at stake changes from acquiring something to avoiding the loss of an option, and loss weighs more heavily than the equivalent gain.

That is what makes the proxy exploitable. Availability is a claim about the world that can be manufactured more cheaply than the condition itself --- a number can be limited, a date can be set, a supply can be described as failing. Nothing about the object has to change. And because the resulting pressure is experienced as an appraisal of value rather than as a reaction to a deadline, the target does not notice that the judgement has been moved.
\end{mechanism}
\begin{tells}
  \item Your estimate of what the thing is worth rose when you learned how little of it there is.
  \item Availability is stated but never evidenced.
  \item The scarcity arrived at the moment a decision was required.
  \item You are competing for access rather than evaluating the thing.
  \item Losing the option feels worse than the thing feels good.
\end{tells}
\begin{moves}
  \item Limit the number, the window or the eligibility, and make the limit visible before the request.
  \item Describe the limit as a fact about the world rather than as a decision you made.
  \item Frame the offer as an option that will disappear rather than as a thing that can be acquired.
  \item Withdraw access gradually so that each stage reads as a real reduction.
\end{moves}
\begin{counters}
  \item Ask what the thing is worth when it is plentiful. If the answer changes with the count, you were never evaluating the thing.
  \item Establish the base rate: how often has this been scarce, and did the scarcity pass without consequence last time?
  \item Separate use from possession. Scarcity inflates the value of owning an option, which is usually not what you actually want.
  \item Check whether the limit is real by trying to exceed it.
  \item Delay past the deadline and observe whether the deadline moves.
\end{counters}
\begin{cost}
Manufactured scarcity is the easiest of the techniques to expose and the most damaging once exposed, because a limit that moves on request was never a limit. It also selects badly: people who respond most strongly to scarcity are frequently the worst customers, since the pressure that brought them produces regret shortly afterwards. And habitual exposure to it degrades the practitioner's own judgement --- someone who always works under artificial deadlines stops being able to tell an urgent thing from a merely announced one.
\end{cost}

\sources{B4}
\seesources
\seealso{T-10-02, T-10-03, T-14-02, T-24-03}
